\chapter{Interactive proofs and PCPs}
\label{chap:interaction}

This chapter covers interactive proof systems ($\cc{IP}$, $\cc{MA}$,
$\cc{AM}$) and probabilistically checkable proofs. Verifiers are machine-free:
a verifier is a polynomial-time function or a language in $\cc{P}$ applied to
an encoded view, and its coins are a uniformly random point of
$\{0,1\}^t$. Formalized so far: the three interactive classes with
completeness $2/3$ and soundness $1/3$, $\cc{P} \subseteq \cc{MA}$ and
$\cc{P} \subseteq \cc{AM}$, $\cc{IP} \subseteq \cc{PSPACE}$ by evaluating the
protocol's game tree, the one-variable multilinear extension, and the PCP
theorem $\cc{NP} = \cc{PCP}(O(\log n), O(1))$ (completeness $1$, soundness
$1/2$) with a constructible randomness bound, proved along Dinur's
gap-amplification route. The main open direction is
$\cc{PSPACE} \subseteq \cc{IP}$ (hence $\cc{IP} = \cc{PSPACE}$) through
sum-check and arithmetization, followed by hardness of approximation built on
the PCP theorem.

\section{Interactive protocols}

\begin{definition}[Protocols, strategies and transcripts]
  \label{def:interaction-protocol}
  \lean{Complexity.Transcript, Complexity.ProverStrategy,
    Complexity.ProverStrategy.Bounded, Complexity.protocolView,
    Complexity.Protocol, Complexity.Protocol.transcript,
    Complexity.Protocol.Accepts, Complexity.Protocol.acceptEvent}
  \leanok
  \uses{def:FP, def:P, def:pair, def:encodings-data,
    def:encodings-bitstring-list}
  A transcript is a list of messages (bit strings); the verifier speaks at
  even positions and the prover at odd ones. A prover strategy is an arbitrary
  function from the visible transcript to the next message; it never sees the
  verifier's coins, and it is bounded by $m$ if the message it returns on
  every transcript has length at most $m$. A protocol consists of a round
  count, a coin count $c$ and a message-length bound (functions of the input
  length), a next-message function in $\cc{FP}$ applied to the encoded view
  $\mathrm{pair}(\mathrm{pair}(x, r), \langle \tau \rangle)$ of the input $x$,
  the private coins $r$ and the transcript $\tau$ (where
  $\langle \tau \rangle$ is the rose-tree serialization of $\tau$), a proof
  that for every $x$, $r$ and $\tau$ the verifier's message has length at most
  the message bound at $|x|$, and a verdict language in $\cc{P}$. Each round
  appends the verifier's message, computed from the current view, and then the
  prover's reply to the transcript extended by that message. The verifier
  accepts if the view of the transcript after the round count at $|x|$ many
  rounds lies in the verdict language, and the accept event is the set of coin
  strings $r \in \{0,1\}^{c(|x|)}$ on which it accepts.
\end{definition}

\begin{definition}[IP]
  \label{def:IP}
  \lean{Complexity.IP}
  \leanok
  \uses{def:interaction-protocol, def:encodings-eventprob}
  $L \in \cc{IP}$ if some protocol whose round count, coin count $c$ and
  message bound $m$ are given by polynomials with natural-number coefficients
  satisfies: for $x \in L$ some strategy bounded by $m(|x|)$ is accepted with
  probability at least $2/3$, and for $x \notin L$ every strategy bounded by
  $m(|x|)$ is accepted with probability at most $1/3$, the probability being
  over uniform $r \in \{0,1\}^{c(|x|)}$. The counts are required to be
  polynomials, not merely polynomially bounded, so that the verifier can
  compute them.
\end{definition}

\begin{definition}[MA and AM]
  \label{def:MA-AM}
  \lean{Complexity.merlinEvent, Complexity.MA, Complexity.arthurEvent,
    Complexity.AM}
  \leanok
  \uses{def:P, def:pair, def:encodings-eventprob}
  $L \in \cc{MA}$ if there are a polynomial $p$ with natural-number
  coefficients and $V \in \cc{P}$ such that
  for $x \in L$ some proof $w$ with $|w| \le p(|x|)$ makes
  $\mathrm{pair}(\mathrm{pair}(x, w), r) \in V$ for at least $2/3$ of the
  $r \in \{0,1\}^{p(|x|)}$, and for $x \notin L$ every such proof does so for at
  most $1/3$ of them. $L \in \cc{AM}$ if the coins $r$ come first and are
  public: the fraction of $r \in \{0,1\}^{p(|x|)}$ admitting some reply $w$
  with $|w| \le p(|x|)$
  and $\mathrm{pair}(\mathrm{pair}(x, r), w) \in V$ is at least $2/3$ for
  members and at most $1/3$ for non-members.
\end{definition}

\begin{proposition}[$\cc{P} \subseteq \cc{MA}$ and $\cc{P} \subseteq \cc{AM}$]
  \label{prop:interaction-P-subset-MA-AM}
  \lean{Complexity.P_subset_MA, Complexity.P_subset_AM}
  \leanok
  \uses{def:P, def:MA-AM}
  $\cc{P} \subseteq \cc{MA}$ and $\cc{P} \subseteq \cc{AM}$.
\end{proposition}
\begin{proof}
  \leanok
  Take $p = 0$, so there are no coins and the proof or reply is empty; the
  verifier decides the input it recovers from the encoded view, so the
  acceptance probability is $1$ on members and $0$ on non-members.
\end{proof}

\begin{lemma}[Strategy extensionality]
  \label{lem:interaction-strategy-ext}
  \lean{Complexity.Protocol.transcript_length,
    Complexity.Protocol.transcript_congr, Complexity.Protocol.accepts_congr,
    Complexity.Protocol.acceptEvent_congr}
  \leanok
  \uses{def:interaction-protocol}
  The transcript after $n$ rounds has length exactly $2n$. If two strategies
  agree on all transcripts of length at most $2n$, then for every input and
  every coin string the transcripts after $n$ rounds coincide. Hence
  acceptance and the accept event on input $x$ depend on the strategy only
  through its values on transcripts of length at most twice the round count
  at $|x|$.
\end{lemma}
\begin{proof}
  \leanok
  Induction on the number of rounds; the transcript after $n$ rounds has
  length exactly $2n$.
\end{proof}

\begin{lemma}[The game value and an optimal prover]
  \label{lem:interaction-game-value}
  \lean{Complexity.Protocol.consFinset, Complexity.Protocol.vset,
    Complexity.Protocol.gval, Complexity.Protocol.card_acceptEvent_le_gval,
    Complexity.Protocol.optStrategy, Complexity.Protocol.optStrategy_bounded,
    Complexity.Protocol.gval_eq_card_acceptEvent}
  \leanok
  \uses{def:interaction-protocol}
  Fix a protocol, an input $x$, a coin count $t$ and a message bound $m$. A
  coin string $r \in \{0,1\}^t$ is consistent with a transcript $\tau$ if
  every verifier message recorded in $\tau$ is the one the verifier sends on
  coins $r$. Define $\mathrm{gval}$ at a transcript $\tau$ by recursion on the
  number of remaining rounds: with none left it counts the coin strings
  consistent with $\tau$ on which the verdict accepts the view of $\tau$, and
  otherwise it sums, over the verifier messages that the consistent coin
  strings produce next, the maximum over replies of length at most $m$ of the
  value one round down. Let $\mathrm{gval}(x)$ be its value at the root: coin
  count $c(|x|)$, message bound and round count at $|x|$, and the empty
  transcript. Then every strategy bounded by the message bound at $|x|$ has at
  most $\mathrm{gval}(x)$ accepting coin strings, and a bounded strategy
  chosen for $x$ (a maximizing reply at every node) has exactly
  $\mathrm{gval}(x)$.
\end{lemma}
\begin{proof}
  \leanok
  Induction on the remaining rounds; the maximizing reply is chosen
  node by node.
\end{proof}

\begin{theorem}[IP membership is a game-value comparison]
  \label{thm:interaction-ip-game-value}
  \lean{Complexity.Protocol.mem_iff_gval, Complexity.IP_membership_is_game_value}
  \leanok
  \uses{def:IP, lem:interaction-game-value}
  If a protocol and a language $L$ satisfy the completeness and soundness
  conditions of $\cc{IP}$ (no polynomial bound on the counts is needed), then
  for every $x$, $x \in L$ if and only if $2^{c(|x|)} < 2\,\mathrm{gval}(x)$,
  where $c$ is the coin count and $\mathrm{gval}(x)$ is the root value of
  Lemma~\ref{lem:interaction-game-value}. In particular, for every
  $L \in \cc{IP}$ there is a protocol with this property.
\end{theorem}
\begin{proof}
  \leanok
  \uses{lem:interaction-game-value}
  By the lemma, the root value is the largest accepting count of a bounded
  strategy. Completeness puts it at least two thirds of the coin space on
  members and soundness at most one third on non-members, so comparing it
  with half the coin space separates them.
\end{proof}

\begin{theorem}[$\cc{IP} \subseteq \cc{PSPACE}$]
  \label{thm:interaction-ip-subset-pspace}
  \lean{Complexity.IP_subset_PSPACE}
  \leanok
  \uses{def:IP, def:PSPACE}
  $\cc{IP} \subseteq \cc{PSPACE}$.
\end{theorem}
\begin{proof}
  \leanok
  \uses{thm:interaction-ip-game-value, lem:space-iterate}
  The game value is written as nested counter loops and walked depth-first on
  a stack with one frame per round; one step of the walk is a polynomial-time
  function, and Lemma~\ref{lem:space-iterate} turns the iteration into a
  polynomial-space machine.
\end{proof}

\begin{proposition}[IP = PSPACE, given the hard half]
  \label{prop:interaction-ip-eq-pspace-of}
  \lean{Complexity.PSPACESubsetIP, Complexity.IP_eq_PSPACE_of_pspace_subset_ip,
    Complexity.IP_eq_PSPACE_of}
  \leanok
  \uses{def:IP, def:PSPACE}
  If $\cc{PSPACE} \subseteq \cc{IP}$, then $\cc{IP} = \cc{PSPACE}$. The
  hypothesis is also recorded, unproved, as the proposition
  \texttt{PSPACESubsetIP}, and the conclusion is stated from either form.
\end{proposition}
\begin{proof}
  \leanok
  \uses{thm:interaction-ip-subset-pspace}
  Antisymmetry with Theorem~\ref{thm:interaction-ip-subset-pspace}.
\end{proof}

\begin{definition}[Completeness and soundness parameters]
  \label{def:interaction-ip-params}
  \uses{def:IP, def:MA-AM}
  $\cc{IP}(c, s)$, $\cc{MA}(c, s)$ and $\cc{AM}(c, s)$: the classes with
  completeness $c$ and soundness $s$ as parameters, together with monotonicity
  in both thresholds and the identification of the hard-wired classes with
  $c = 2/3$, $s = 1/3$.
\end{definition}

\begin{theorem}[Sequential repetition]
  \label{thm:interaction-sequential-repetition}
  \uses{def:interaction-ip-params}
  For constants $0 \le s < c \le 1$ and every polynomial $q$,
  $\cc{IP}(c, s) = \cc{IP}(1 - 2^{-q(n)}, 2^{-q(n)})$: run the protocol
  sequentially polynomially many times and accept by a threshold vote.
\end{theorem}
\begin{proof}
  \uses{lem:interaction-strategy-ext}
  Soundness against an adaptive prover needs the conditional per-run bound
  across rounds, not only independence of the coins; strategy extensionality
  lets each run's strategy be read off the combined one.
\end{proof}

\begin{proposition}[$\cc{NP} \subseteq \cc{MA}$]
  \label{prop:interaction-np-subset-ma}
  \uses{def:NP, def:MA-AM}
  $\cc{NP} \subseteq \cc{MA}$.
\end{proposition}
\begin{proof}
  Merlin sends the NP witness and Arthur ignores his coins.
\end{proof}

\begin{proposition}[$\cc{MA} \subseteq \cc{IP}$]
  \label{prop:interaction-ma-subset-ip}
  \uses{def:MA-AM, def:IP}
  $\cc{MA} \subseteq \cc{IP}$.
\end{proposition}
\begin{proof}
  A one-round protocol with an empty verifier message; the prover's reply is
  Merlin's proof, which cannot depend on the private coins.
\end{proof}

\begin{corollary}[$\cc{NP} \subseteq \cc{IP}$]
  \label{cor:interaction-np-subset-ip}
  \uses{def:NP, def:IP}
  $\cc{NP} \subseteq \cc{IP}$.
\end{corollary}
\begin{proof}
  \uses{prop:interaction-np-subset-ma, prop:interaction-ma-subset-ip}
  Compose the two inclusions.
\end{proof}

\section{Algebraic tools and IP = PSPACE}

\begin{lemma}[One-variable multilinear extension]
  \label{lem:interaction-mle1}
  \lean{Complexity.mle₁, Complexity.mle₁_zero, Complexity.mle₁_one,
    Complexity.mle₁_bool}
  \leanok
  For $f : \{0,1\} \to R$ with $R$ a commutative ring, let
  $\widetilde{f} : R \to R$ be $\widetilde{f}(x) = f(0)(1 - x) + f(1)\,x$.
  Then $\widetilde{f}(0) = f(0)$ and
  $\widetilde{f}(1) = f(1)$, so $\widetilde{f}$ agrees with $f$ at both Boolean
  points.
\end{lemma}
\begin{proof}
  \leanok
  Direct computation.
\end{proof}

\begin{definition}[Multilinear extension]
  \label{def:interaction-mle}
  \uses{lem:interaction-mle1}
  For $f : \{0,1\}^m \to \mathbb{F}$, the multilinear extension
  $\widetilde{f}(x) = \sum_{b \in \{0,1\}^m} f(b) \prod_i
  (b_i x_i + (1 - b_i)(1 - x_i))$, with the facts that it agrees with $f$ on the
  cube and is the unique multilinear polynomial doing so.
\end{definition}

\begin{lemma}[Schwartz--Zippel over coin strings]
  \label{lem:interaction-schwartz-zippel}
  If $p$ is a nonzero polynomial over a finite field $\mathbb{F}$ of total
  degree at most $d$ in $m$ variables, a uniformly random point of
  $\mathbb{F}^m$ is a root with probability at most $d / |\mathbb{F}|$; in
  particular a nonzero univariate polynomial of degree $d$ has at most $d$
  roots. Mathlib provides the counting statement for
  \texttt{MvPolynomial} (\texttt{MvPolynomial.schwartz\_zippel\_totalDegree},
  over an integral domain with points drawn from $S^m$ for a finite set $S$);
  nothing in this node is formalized in the library yet, and what remains is
  the bridge to the library's finite probability over encoded coin strings.
\end{lemma}

\begin{definition}[The sum-check protocol]
  \label{def:interaction-sumcheck}
  \uses{def:interaction-protocol}
  For a polynomial $g$ in $m$ variables over $\mathbb{F}$ with individual
  degree at most $d$, given by an evaluation oracle, and a claimed value $v$ of
  $\sum_{b \in \{0,1\}^m} g(b)$: in round $i$ the prover sends a univariate
  polynomial $h_i$ of degree at most $d$, the verifier checks
  $h_i(0) + h_i(1)$ against the current claim, and replaces the claim by
  $h_i(r_i)$ for a random $r_i \in \mathbb{F}$; at the end it checks
  $g(r_1, \dots, r_m)$ against the last claim.
\end{definition}

\begin{theorem}[Sum-check completeness and soundness]
  \label{thm:interaction-sumcheck}
  \uses{def:interaction-sumcheck, lem:interaction-schwartz-zippel}
  \notready
  If the claim is correct, the honest prover is accepted with probability $1$.
  If it is incorrect, every prover is accepted with probability at most
  $m d / |\mathbb{F}|$.
\end{theorem}
\begin{proof}
  \uses{lem:interaction-schwartz-zippel}
  Round-by-round invariant: a false claim stays false unless the verifier hits
  a root of the difference between the prover's polynomial and the true one.
\end{proof}

\begin{definition}[Arithmetization of QBF with degree reduction]
  \label{def:interaction-arithmetization}
  \uses{def:space-qbf, def:interaction-mle}
  Map a QBF to an expression over $\mathbb{F}$: $\neg a \mapsto 1 - a$,
  $a \wedge b \mapsto ab$, $\forall x\, \varphi \mapsto \varphi|_{x=0} \cdot
  \varphi|_{x=1}$, $\exists x\, \varphi \mapsto 1 - (1 - \varphi|_{x=0})(1 -
  \varphi|_{x=1})$, interleaved with linearization operators
  $L_{x}\, p = (1 - x)\, p|_{x=0} + x\, p|_{x=1}$ that keep every intermediate
  polynomial of individual degree at most $2$.
\end{definition}

\begin{lemma}[Arithmetization is faithful and degree-controlled]
  \label{lem:interaction-degree-reduction}
  \uses{def:interaction-arithmetization, lem:space-qbf-semantics}
  \notready
  On Boolean points the arithmetized expression takes the value $1$ exactly on
  true formulas and $0$ otherwise, linearization does not change values on
  Boolean points, and with linearization every polynomial the protocol handles
  has individual degree at most $2$ and total size polynomial in the formula.
\end{lemma}

\begin{theorem}[An interactive proof for TQBF]
  \label{thm:interaction-tqbf-protocol}
  \uses{def:space-tqbf, def:IP}
  \notready
  $\cc{TQBF} \in \cc{IP}$, with polynomial verifier time, polynomial
  communication, perfect completeness and soundness error bounded away from
  $1$, over a field whose size dominates the degree and soundness parameters
  and admits a polynomial-size encoding.
\end{theorem}
\begin{proof}
  \uses{thm:interaction-sumcheck, lem:interaction-degree-reduction,
    thm:interaction-sequential-repetition}
  Run a sum-check-style round for each quantifier and linearization operator,
  from the outside in. Naive arithmetization without degree reduction has
  exponential degree, which is why the degree-control lemma is a separate
  node.
\end{proof}

\begin{theorem}[$\cc{PSPACE} \subseteq \cc{IP}$ (Shamir)]
  \label{thm:interaction-pspace-subset-ip}
  \uses{def:PSPACE, def:IP}
  $\cc{PSPACE} \subseteq \cc{IP}$. The statement is recorded, unproved, as the
  proposition \texttt{PSPACESubsetIP} of
  Proposition~\ref{prop:interaction-ip-eq-pspace-of}.
\end{theorem}
\begin{proof}
  \uses{thm:interaction-tqbf-protocol, thm:space-tqbf-complete}
  Compose a polynomial-time reduction to TQBF with the TQBF protocol; $\cc{IP}$
  is closed under polynomial-time many-one reductions.
\end{proof}

\begin{theorem}[$\cc{IP} = \cc{PSPACE}$]
  \label{thm:interaction-ip-eq-pspace}
  \uses{def:IP, def:PSPACE}
  $\cc{IP} = \cc{PSPACE}$.
\end{theorem}
\begin{proof}
  \uses{thm:interaction-pspace-subset-ip, prop:interaction-ip-eq-pspace-of}
  Apply Proposition~\ref{prop:interaction-ip-eq-pspace-of}.
\end{proof}

\section{Probabilistically checkable proofs}

\begin{definition}[PCP verifiers]
  \label{def:interaction-pcp-verifier}
  \lean{Complexity.PCPVerifier, Complexity.PCPVerifier.answers,
    Complexity.PCPVerifier.Accepts, Complexity.PCPVerifier.acceptEvent,
    Complexity.PCPVerifier.QueryBounded}
  \leanok
  \uses{def:FP, def:P, def:pair, def:encodings-data}
  A PCP verifier is non-adaptive: it has a list of proof positions (natural
  numbers) for each input $x$ and coin string $r$, together with an $\cc{FP}$
  function mapping $\mathrm{pair}(x, r)$ to the rose-tree serialization of
  that list, and a verdict language in $\cc{P}$ that decides
  $\mathrm{pair}(\mathrm{pair}(x, r), a)$, where $a$ lists the bits of the
  proof (a finite bit string) found at those positions. A position past the
  end of the proof reads as $0$. The verifier is query-bounded by $q$ if it
  reads at most $q(|x|)$ positions for every $x$ and every coin string $r$.
\end{definition}

\begin{definition}[PCP classes]
  \label{def:PCP}
  \lean{Complexity.PCP}
  \leanok
  \uses{def:interaction-pcp-verifier, def:encodings-eventprob}
  For arbitrary functions $r, q : \mathbb{N} \to \mathbb{N}$,
  $L \in \cc{PCP}(r, q)$ if some verifier query-bounded by $q$, using exactly
  $r(|x|)$ uniformly random coins, accepts some proof of each $x \in L$ with
  probability $1$, and accepts every proof of each $x \notin L$ with
  probability at most $1/2$.
\end{definition}

\begin{definition}[Constructible bounds]
  \label{def:interaction-constructible}
  \lean{Complexity.Constructible}
  \leanok
  \uses{def:FP}
  A function $r$ is constructible if $x \mapsto 1^{r(|x|)}$ is in $\cc{FP}$.
\end{definition}

\begin{proposition}[Elementary PCP facts]
  \label{prop:interaction-P-subset-PCP}
  \lean{Complexity.P_subset_PCP, Complexity.PCP_mono_queries}
  \leanok
  \uses{def:P, def:PCP}
  $\cc{P} \subseteq \cc{PCP}(r, q)$ for all $r$ and $q$, and
  $\cc{PCP}(r, q) \subseteq \cc{PCP}(r, q')$ whenever $q \le q'$ pointwise.
\end{proposition}
\begin{proof}
  \leanok
  For the first, a verifier that reads nothing, ignores its coins and decides
  the input it recovers from the encoded view. For the second, the same
  verifier meets any pointwise larger query bound. Monotonicity in the coin
  count is not formalized.
\end{proof}

\begin{proposition}[The randomness bound must be constructible]
  \label{prop:interaction-constructible-needed}
  \lean{Complexity.lengthLang_mem_iUnion_PCP}
  \leanok
  \uses{def:PCP}
  For every set $A \subseteq \mathbb{N}$, computable or not, the language
  $\{x : |x| \in A\}$ lies in $\bigcup \cc{PCP}(r, q)$ over all
  $r = O(\log n)$ and $q = O(1)$, with no constructibility requirement on $r$.
\end{proposition}
\begin{proof}
  \leanok
  Use the indicator of $A$ as the coin count, the zero query bound, and a
  verifier that reads nothing and accepts exactly when it received one coin.
  Since there are uncountably many sets $A$
  and $\cc{NP}$ is countable, the unrestricted union is not $\cc{NP}$; this
  consequence is argued in the documentation, not stated in Lean.
\end{proof}

\begin{lemma}[PCP soundness amplification]
  \label{lem:interaction-pcp-amplification}
  \lean{Complexity.PCPWith, Complexity.PCPWith_half, Complexity.PCPWith_square,
    Complexity.mem_PCP_of_PCPWith}
  \leanok
  \uses{def:PCP, def:interaction-constructible}
  For a rational $s$, let $\cc{PCP}_s(r, q)$ be the class with soundness
  error $s$ in place of $1/2$, so $\cc{PCP}_{1/2}(r, q) = \cc{PCP}(r, q)$. If
  $r$ is constructible and $0 \le s$, then
  $\cc{PCP}_s(r, q) \subseteq \cc{PCP}_{s^2}(2r, 2q)$. Consequently, if $r$
  is constructible and $0 \le s < 1$, every language in $\cc{PCP}_s(r, q)$ is
  in $\cc{PCP}(2^j r, 2^j q)$ for some $j \in \mathbb{N}$.
\end{lemma}
\begin{proof}
  \leanok
  Running the verifier twice on independent coins squares the error and keeps
  perfect completeness; the coin string is split at the constructible per-run
  count, and doubling keeps the bound constructible. After $j$ doublings the
  error is $s^{2^j} \le s^j$, which is below $1/2$ for large $j$.
\end{proof}

\begin{definition}[Constraint graphs]
  \label{def:interaction-constraint-graph}
  \lean{Complexity.ConstraintGraph, Complexity.ConstraintGraph.Satisfiable,
    Complexity.ConstraintGraph.unsatFrac, Complexity.ConstraintGraph.unsatVal,
    Complexity.ConstraintGraph.unsatVal_eq_zero_iff_satisfiable}
  \leanok
  A constraint graph over an alphabet $\Sigma$ has vertices
  $\{0, \dots, V - 1\}$ and finitely many indexed directed edges (parallel
  edges and self-loops allowed), each carrying a Boolean-valued binary
  constraint on the labels of its tail and head. It is satisfiable if some
  labelling of the vertices satisfies every edge. For finite nonempty
  $\Sigma$, its unsatisfiability value is the least, over all labellings, of
  the fraction of edges left unsatisfied (taken to be $0$ for an edgeless
  graph); it is $0$ exactly when the graph is satisfiable.
\end{definition}

\begin{theorem}[Dinur's gap theorem for 3SAT]
  \label{thm:interaction-dinur-gap}
  \lean{Complexity.DinurAlpha, Complexity.Amplifier,
    Complexity.Amplifier.dichotomy, Complexity.dinurAmp, Complexity.gapGraph,
    Complexity.satisfiable_gapGraph, Complexity.gap_le_unsatVal_gapGraph,
    Complexity.numEdges_gapGraph_le, Complexity.dinurAmp_gap_pos,
    Complexity.dinurAmp_gap_le_one}
  \leanok
  \uses{def:interaction-constraint-graph, def:polytime-sat}
  There are a fixed finite alphabet, constants $0 < \gamma \le 1$ and $c$,
  and a map $\varphi \mapsto G_\varphi$ from CNF formulas to constraint graphs
  over that alphabet such that: for every CNF $\varphi$ with $m$ clauses,
  $G_\varphi$ has at most $(6m + 2)^{c} \cdot 3m$ edges; and for every
  $\varphi$ in which every clause has exactly three literals, $G_\varphi$ is
  satisfiable if $\varphi$ is, and the unsatisfiability value of $G_\varphi$
  is at least $\gamma$ if $\varphi$ is unsatisfiable. The map is defined as a
  noncomputable mathematical reduction; the polynomial-time counterpart used
  for the PCP theorem is built from the same amplifier inside the proof of
  Theorem~\ref{thm:interaction-np-subset-pcp}.
\end{theorem}
\begin{proof}
  \leanok
  An amplifier is a transformation of constraint graphs over a fixed alphabet
  that multiplies the edge count by at most a constant, preserves
  satisfiability, and maps unsatisfiability value $v$ to at least
  $\min(\gamma, 2v)$. Dinur's round (expander preprocessing, graph powering,
  and alphabet reduction through a Hadamard-code assignment tester, over an
  explicit expander family) is one. Iterating it $k$ times on a graph with at
  most $2^k$ edges gives the dichotomy: satisfiable graphs stay satisfiable
  and unsatisfiable ones reach value at least $\gamma$. Apply it to the
  $3m$-edge constraint graph of the formula with
  $k = \lfloor \log_2 3m \rfloor + 1$.
\end{proof}

\begin{theorem}[$\cc{NP} \subseteq \cc{PCP}(O(\log n), O(1))$]
  \label{thm:interaction-np-subset-pcp}
  \lean{Complexity.exists_pcp_of_mem_NP}
  \leanok
  \uses{def:NP, def:PCP, def:interaction-constructible}
  For every $L \in \cc{NP}$ there are a constructible $r$ with
  $r = O(\log n)$ and a $q$ with $q = O(1)$ such that
  $L \in \cc{PCP}(r, q)$. Here $f = O(g)$ means $f(n) \le C\,g(n)$ for some
  constant $C$ and all sufficiently large $n$, and $\log n$ is
  $\lfloor \log_2 n \rfloor$; the $q$ produced is a constant function.
\end{theorem}
\begin{proof}
  \leanok
  \uses{thm:cook-levin, lem:polytime-tseitin, def:interaction-constraint-graph,
    thm:interaction-dinur-gap, lem:interaction-pcp-amplification}
  Map each input, by an $\cc{FP}$ function, to an exact-3CNF that is
  satisfiable exactly when the input is in $L$ (Cook--Levin followed by
  Tseitin splitting). Pad its constraint graph to a size fixed by the input
  length and compute the gap constraint graph by a polynomial-time
  implementation of every round of Dinur's amplifier (the amplifier of
  Theorem~\ref{thm:interaction-dinur-gap}, over its explicit expander
  family). The verifier picks a random edge from logarithmically many coins
  and reads the fixed-width labels of its two endpoints, so members are
  accepted with probability $1$ and non-members with probability bounded
  away from $1$; Lemma~\ref{lem:interaction-pcp-amplification} brings the
  soundness error below $1/2$.
\end{proof}

\begin{theorem}[$\cc{PCP}(O(\log n), O(1)) \subseteq \cc{NP}$]
  \label{thm:interaction-pcp-subset-np}
  \lean{Complexity.PCP_subset_NP}
  \leanok
  \uses{def:PCP, def:NP, def:interaction-constructible}
  If $r$ is constructible, $r = O(\log n)$ and $q = O(1)$ (in the sense of
  Theorem~\ref{thm:interaction-np-subset-pcp}), then
  $\cc{PCP}(r, q) \subseteq \cc{NP}$.
\end{theorem}
\begin{proof}
  \leanok
  \uses{lem:polytime-guess-verify}
  The verifier's behaviour on all $2^{r(n)} = \mathrm{poly}(n)$ coin strings
  is determined by a polynomial-size table listing, for each coin string, the
  answers to its constantly many queries. The witness is such a table; a
  polynomial-time check, which uses constructibility to write out
  $2^{r(n)}$, verifies its length, its consistency (a position queried under
  two coin strings gets the same answer), and acceptance on every coin
  string. A member has such a table by perfect completeness; a non-member has
  none, since a consistent table gives a proof accepted with probability
  $1 > 1/2$.
\end{proof}

\begin{theorem}[PCP theorem]
  \label{thm:interaction-pcp-theorem}
  \lean{Complexity.PCP_theorem}
  \leanok
  \uses{def:NP, def:PCP, def:interaction-constructible}
  $\cc{NP} = \bigcup \cc{PCP}(r, q)$, the union over all constructible $r$
  with $r = O(\log n)$ and all $q$ with $q = O(1)$ (in the sense of
  Theorem~\ref{thm:interaction-np-subset-pcp}; no constructibility is
  required of $q$).
\end{theorem}
\begin{proof}
  \leanok
  \uses{thm:interaction-np-subset-pcp, thm:interaction-pcp-subset-np}
  The two inclusions. The constructibility requirement on $r$ cannot be
  dropped, by
  Proposition~\ref{prop:interaction-constructible-needed}.
\end{proof}

\section{Hardness of approximation}

\begin{definition}[Gap constraint satisfaction problems]
  \label{def:interaction-gap-csp}
  \uses{def:promise, def:interaction-constraint-graph}
  For $0 \le s < c \le 1$ and fixed arity and alphabet, $\mathrm{Gap}_{c,s}$-CSP
  is the promise problem whose yes-instances are encoded constraint systems
  with value at least $c$ and whose no-instances have value less than $s$.
  $\mathrm{Gap}$-MAX-3SAT$_{1,s}$ separates satisfiable 3CNFs from those in
  which every assignment satisfies fewer than an $s$ fraction of the clauses.
\end{definition}

\begin{theorem}[Gap MAX-3SAT is NP-hard]
  \label{thm:interaction-gap-max3sat}
  \uses{def:interaction-gap-csp, def:NP}
  \notready
  There is $\varepsilon > 0$ such that $\mathrm{Gap}$-MAX-3SAT$_{1, 1 -
  \varepsilon}$ is NP-hard under polynomial-time reductions to the promise
  problem.
\end{theorem}
\begin{proof}
  \uses{thm:interaction-pcp-theorem, thm:interaction-np-subset-pcp}
  Turn each random string of the PCP verifier into a constant-size CNF over the
  proof bits, then convert to 3CNF; the constant soundness gap becomes a
  constant fraction of violated clauses. The equivalence between PCP
  characterizations and gap-CSP hardness is the reusable statement.
\end{proof}

\begin{definition}[Label cover]
  \label{def:interaction-label-cover}
  \uses{def:interaction-constraint-graph}
  A label-cover instance is a bipartite constraint graph with label sets
  $\Sigma_A$, $\Sigma_B$ whose constraints are projections
  $\pi_e : \Sigma_A \to \Sigma_B$; its value is the largest fraction of edges
  satisfied by a labelling.
\end{definition}

\begin{theorem}[Gap label cover is NP-hard]
  \label{thm:interaction-label-cover-hard}
  \uses{def:interaction-label-cover, def:interaction-gap-csp}
  \notready
  For some constant $\delta < 1$ and constant-size label sets, distinguishing
  label-cover instances of value $1$ from those of value at most $\delta$ is
  NP-hard.
\end{theorem}
\begin{proof}
  \uses{thm:interaction-gap-max3sat}
  The clause--variable game on a gap-3SAT instance.
\end{proof}

\begin{theorem}[Parallel repetition]
  \label{thm:interaction-parallel-repetition}
  \uses{def:interaction-label-cover}
  \notready
  For every label-cover instance of value at most $\delta < 1$, its $k$-fold
  parallel repetition has value at most $\delta'^{\,k}$ for a constant
  $\delta' < 1$ depending only on $\delta$ and the label-set sizes. Hence
  gap label cover with arbitrarily small constant soundness is NP-hard.
\end{theorem}
\begin{proof}
  \uses{thm:interaction-label-cover-hard}
  Raz's theorem, or a simpler projection-game variant with a weaker exponent.
\end{proof}
